\documentclass[11pt,a4paper]{article}
\usepackage{auditstyle}
\usepackage{bm}

\newcommand{\eps}{\varepsilon}
\newcommand{\W}{\mathcal W}
\newcommand{\T}{\mathbb T}
\newcommand{\dd}{\,d}
\newtheorem{corollary}[theorem]{Corollary}

\title{Part II free fixed-coarse rate\\[3pt]
\large An explicit \(D4\) filter estimate and covariance-symbol proof}
\author{Lamport-form proof}
\date{9 August 2026}

\begin{document}
\maketitle

\begin{resultbox}
\textbf{Quantitative free theorem.}  Fix a coarse scale \(M\), a torus
\(\T_L\), and a coarse Weyl generator \(W_M(\eps_Mq+ip)\).  For the MMST
nearest-neighbour free vacua and the exact Daubechies \(D4\) observable
embeddings,
\[
 \left|
 \omega_{M+r,L}^{(0)}\!\left(\alpha_M^{M+r}
      (W_M(\eps_Mq+ip))\right)
 -\omega_L^{(0)}\!\left(\beta_M(W_M(\eps_Mq+ip))\right)
 \right|
 \le C_{M,L,m,q,p}\,2^{-\theta r},                              \tag{R}
\]
for all sufficiently large \(r\), where
\[
 \eta=2(1-\log_2\sqrt3)>0,\qquad
 \theta=\frac{2\eta}{5+\eta}>0.                                \tag{E}
\]
Every constant is displayed in Section \ref{sec:rate}.  The theorem holds in
the correct fixed-coarse topology and makes no diagonal \(M=N\) assertion.
\end{resultbox}

\tableofcontents

\section{Exact covariance formulas}\label{sec:cov}

Fix \(M\), and put \(N=M+r\).  Use all conventions from the companion model
sheet.  Write
\[
 \Gamma_\infty:=\frac{\pi}{L}\mathbb Z,\qquad
 \gamma(k):=\sqrt{m^2+k^2},\qquad
 \gamma_N(k):=\sqrt{m^2+4\eps_N^{-2}\sin^2(\eps_Nk/2)}.          \tag{1.1}
\]
Extend the coarse Fourier transforms \(\widehat q_M,\widehat p_M\) periodically
from \(\Gamma_M\) to \(\Gamma_\infty\), with period \(2\pi/\eps_M\).

Define the finite and infinite mask products
\[
 P_r(\ell):=\prod_{j=1}^r m_0(2^{-j}\ell),\qquad
 P_\infty(\ell):=\prod_{j=1}^\infty m_0(2^{-j}\ell)
 =\widehat\phi(\ell).                                          \tag{1.2}
\]
The equality in (1.2) is the Fourier form of the \(D4\) scaling equation with
\(\widehat\phi(0)=1\).

\begin{lemma}[pulled-back fine covariance]\label{lem:finecov}
For \(\xi=\eps_Mq+ip\),
\[
 \omega_{M+r,L}^{(0)}(\alpha_M^{M+r}(W_M(\xi)))
 =e^{-Q_{r,M}(q,p)/4},                                        \tag{1.3}
\]
where
\[
 Q_{r,M}(q,p)
 =\frac{\eps_M}{2L}\sum_{k\in\Gamma_\infty}
 \mathbf1_{\Gamma_{M+r}}(k)|P_r(\eps_Mk)|^2
 \left(\gamma_{M+r}(k)^{-1}|\widehat q_M(k)|^2
       +\gamma_{M+r}(k)|\widehat p_M(k)|^2\right).             \tag{1.4}
\]
\end{lemma}

\subsection*{Proof in Lamport form}

\LS{1}{1}{Iterate the momentum-space refinement map.}

The one-step MMST map is multiplication by
\(\sqrt2\,m_0(\eps_{j+1}k)\), together with periodic extension of the
coarse Fourier vector.  Therefore
\[
 \widehat{R_M^{M+r}(q,p)}(k)
 =2^{r/2}P_r(\eps_Mk)
   (\widehat q_M(k),\widehat p_M(k)),
 \qquad k\in\Gamma_{M+r}.                                      \tag{1.5}
\]

\LS{1}{2}{Insert (1.5) in the fine vacuum covariance.}

The fine normalized momentum sum has weight
\[
 \frac1{2r_{M+r}}=\frac{2^{-r}}{2r_M}
 =2^{-r}\frac{\eps_M}{2L}.                                    \tag{1.6}
\]
The squared factor \(2^{r/2}\) in (1.5) cancels the \(2^{-r}\) in
(1.6).  Substitution in the exact free state formula gives (1.3)--(1.4).
\QedStep

\begin{lemma}[continuum pullback covariance]\label{lem:ctcov}
\[
 \omega_L^{(0)}(\beta_M(W_M(\xi)))=e^{-Q_{\infty,M}(q,p)/4},   \tag{1.7}
\]
where
\[
 Q_{\infty,M}(q,p)
 =\frac{\eps_M}{2L}\sum_{k\in\Gamma_\infty}
 |P_\infty(\eps_Mk)|^2
 \left(\gamma(k)^{-1}|\widehat q_M(k)|^2
       +\gamma(k)|\widehat p_M(k)|^2\right).                   \tag{1.8}
\]
\end{lemma}

\begin{proof}
The continuum embedding sends \(q,p\) to linear combinations of the periodized
scaling functions.  Their Fourier coefficients equal
\(\eps_M^{1/2}e^{-ikx}\widehat\phi(\eps_Mk)\).  Parseval's identity on
\(\T_L\), followed by the continuum free-vacuum formula, gives (1.8).
Equation (1.2) then identifies the multiplier with \(P_\infty\).
\end{proof}

\section{A uniform algebraic envelope for the exact D4 product}

Recall the exact factorization
\[
 m_0(t)=\left(\frac{1+e^{-it}}2\right)^2D(t),\qquad
 D(t)=\frac{(1+\sqrt3)+(1-\sqrt3)e^{-it}}2,                    \tag{2.1}
\]
and
\[
 |D(t)|^2=2-\cos t.                                            \tag{2.2}
\]
Set
\[
 a:=\log_2\sqrt3<1,\qquad
 \eta:=2-2a=2(1-\log_2\sqrt3)>0,\qquad
 C_D:=48\pi^4e^{1/6}.                                         \tag{2.3}
\]

\begin{lemma}[uniform finite-product decay]\label{lem:envelope}
For every \(r\ge0\) and every \(\ell\) satisfying
\(|\ell|\le\pi2^r\),
\[
 |P_r(\ell)|^2\le C_D(1+|\ell|)^{-2-\eta}.                     \tag{2.4}
\]
The same estimate holds for \(P_\infty\) on \(\mathbb R\).
\end{lemma}

\subsection*{Proof in Lamport form}

\LS{1}{1}{Bound the \(D\)-factor product.}

Let
\[
 n(\ell):=\left\lceil\log_2\max(1,|\ell|)\right\rceil.          \tag{2.5}
\]
By (2.2), \(|D(t)|\le\sqrt3\).  Also
\[
 |D(t)|^2=1+(1-\cos t)\le e^{t^2/2},
\qquad |D(t)|\le e^{t^2/4}.                                   \tag{2.6}
\]
For \(j>n(\ell)\), use (2.6); for \(j\le n(\ell)\), use
\(|D|\le\sqrt3\).  Since
\[
 \sum_{j=n+1}^\infty4^{-j}=\frac{4^{-n}}3,\qquad
 |\ell|^24^{-n(\ell)}\le1,                                    \tag{2.7}
\]
we obtain, for every finite \(r\),
\begin{align}
 \prod_{j=1}^r|D(2^{-j}\ell)|
 &\le(\sqrt3)^{n(\ell)}
 \exp\left(\frac{\ell^2}{4}\sum_{j=n(\ell)+1}^\infty4^{-j}\right)\\
 &\le\sqrt3e^{1/12}(1+|\ell|)^a.                              \tag{2.8}
\end{align}
If \(r<n(\ell)\), the omitted tail factor in the first line is absent and the
same upper bound remains valid.

\LS{1}{2}{Bound the cosine product.}

For \(r\ge1\), the finite sine-product identity gives
\[
 \prod_{j=1}^r\left|\cos\left(\frac{\ell}{2^{j+1}}\right)\right|
 =\frac{|\sin(\ell/2)|}{2^r|\sin(\ell/2^{r+1})|}.              \tag{2.9}
\]
When \(|\ell|\le\pi2^r\), the denominator angle has absolute value at most
\(\pi/2\).  The elementary bound
\(|\sin u|\ge2|u|/\pi\) on that interval yields
\[
 \prod_{j=1}^r\left|\cos\left(\frac{\ell}{2^{j+1}}\right)\right|
 \le\min\left(1,\frac{\pi|\sin(\ell/2)|}{|\ell|}\right).        \tag{2.10}
\]

\LS{1}{3}{Combine the factors.}

Equations (2.1), (2.8), and (2.10) imply
\[
 |P_r(\ell)|^2
 \le3e^{1/6}(1+|\ell|)^{2a}
 \min(1,\pi^4|\ell|^{-4}).                                    \tag{2.11}
\]
For \(|\ell|\le1\), multiply the desired right side by
\((1+|\ell|)^4\le16\).  For \(|\ell|\ge1\), use
\((1+|\ell|)^4|\ell|^{-4}\le16\).  In both cases,
\[
 |P_r(\ell)|^2
 \le48\pi^4e^{1/6}(1+|\ell|)^{-4+2a}
 =C_D(1+|\ell|)^{-2-\eta}.                                   \tag{2.12}
\]
The case \(r=0\), restricted to \(|\ell|\le\pi\), is covered by the same
constant.  Finally, take \(r\to\infty\) in (2.12) and use (1.2) to obtain the
infinite-product bound. \QedStep

\begin{lemma}[explicit filter-tail error]\label{lem:filtertail}
If \(|\ell|\le2^r\), then
\[
 0\le |P_r(\ell)|^2-|P_\infty(\ell)|^2
 \le\frac{\sqrt3}{12}\,\ell^2\,4^{-r}.                         \tag{2.13}
\]
\end{lemma}

\subsection*{Proof in Lamport form}

\LS{1}{1}{Compute the second derivative exactly.}

Let \(u(t):=|m_0(t)|^2\).  From (2.1)--(2.2),
\[
 u(t)=\frac{2+3\cos t-\cos^3t}{4},\quad
 u'(t)=-\frac34\sin^3t,\quad
 u''(t)=-\frac94\sin^2t\cos t.                                \tag{2.14}
\]
The maximum of \(x(1-x^2)\) on \([0,1]\) is
\(2/(3\sqrt3)\), attained at \(x=1/\sqrt3\).  Hence
\[
 \sup_t|u''(t)|=\frac{\sqrt3}{2}.                              \tag{2.15}
\]
Because \(u(0)=1\) and \(u'(0)=0\), Taylor's theorem gives
\[
 0\le1-u(t)\le\frac{\sqrt3}{4}t^2.                             \tag{2.16}
\]
The left inequality follows independently from the QMF identity
\(u(t)+u(t+\pi)=1\).

\LS{1}{2}{Estimate the omitted infinite product.}

Since every \(u(t)\in[0,1]\),
\[
 1-\prod_{j=r+1}^\infty u(2^{-j}\ell)
 \le\sum_{j=r+1}^\infty(1-u(2^{-j}\ell)).                      \tag{2.17}
\]
For \(|\ell|\le2^r\), all arguments on the right have absolute value at most
\(1/2\).  Equations (2.16)--(2.17) give
\[
 1-\prod_{j=r+1}^\infty u(2^{-j}\ell)
 \le\frac{\sqrt3}{4}\ell^2\sum_{j=r+1}^\infty4^{-j}
 =\frac{\sqrt3}{12}\ell^24^{-r}.                              \tag{2.18}
\]
Multiplication by \(|P_r(\ell)|^2\le1\) proves (2.13).
\QedStep

\section{Dispersion error}

\begin{lemma}[nearest-neighbour versus continuum symbol]\label{lem:dispersion}
If \(|\eps_Nk|\le\pi\), then
\begin{align}
 0&\le\gamma(k)-\gamma_N(k)
 \le\frac{\eps_N^2k^4}{24m},                                  \tag{3.1}\\
 0&\le\gamma_N(k)^{-1}-\gamma(k)^{-1}
 \le\frac{\eps_N^2k^4}{24m^3}.                                \tag{3.2}
\end{align}
\end{lemma}

\subsection*{Proof in Lamport form}

\LS{1}{1}{Bound the discrete Laplacian symbol.}

For \(0\le x\le\pi/2\),
\[
 \sin x\le x,\qquad \sin x\ge x-\frac{x^3}{6}.                 \tag{3.3}
\]
For the second inequality, the derivative of
\(\sin x-x+x^3/6\) is nonnegative because
\(\cos x\ge1-x^2/2\).  Since \(x-x^3/6\ge0\) on this interval,
\[
 x^2-\frac{x^4}{3}
 \le\left(x-\frac{x^3}{6}\right)^2
 \le\sin^2x\le x^2.                                           \tag{3.4}
\]
Putting \(x=|\eps_Nk|/2\) gives
\[
 0\le k^2-4\eps_N^{-2}\sin^2(\eps_Nk/2)
 \le\frac{\eps_N^2k^4}{12}.                                   \tag{3.5}
\]

\LS{1}{2}{Pass to square roots and inverses.}

Since both dispersions are at least \(m\),
\[
 \gamma-\gamma_N
 =\frac{\gamma^2-\gamma_N^2}{\gamma+\gamma_N}
 \le\frac{\eps_N^2k^4}{24m},                                  \tag{3.6}
\]
which is (3.1).  Moreover,
\[
 \gamma_N^{-1}-\gamma^{-1}
 =\frac{\gamma-\gamma_N}{\gamma\gamma_N}
 \le\frac{\eps_N^2k^4}{24m^3},                                \tag{3.7}
\]
which is (3.2). \QedStep

\section{The explicit rate}\label{sec:rate}

Put
\[
 A_q:=\sup_{k\in\Gamma_\infty}|\widehat q_M(k)|,\qquad
 A_p:=\sup_{k\in\Gamma_\infty}|\widehat p_M(k)|,                \tag{4.1}
\]
which are finite because the extensions are periodic.  Define
\[
 \kappa:=\frac{\pi}{L},\qquad b:=\eps_M\kappa=\frac{\pi}{r_M},
 \qquad c_0:=\frac{\sqrt3}{12}.                                \tag{4.2}
\]
Set
\begin{align}
 C_{\rm low}
 :=\frac{3\eps_M^3}{2L}\Bigg\{&
 A_q^2\left(\frac{\kappa^4}{24m^3}+\frac{c_0\kappa^2}{m}\right)\\
 &+A_p^2\left(\frac{\kappa^4}{24m}
       +c_0(m+\kappa)\kappa^2\right)\Bigg\},                   \tag{4.3}\\
 C_*&:=2C_D\left(\frac{A_q^2}{m}
              +A_p^2(m+\eps_M^{-1})\right).                   \tag{4.4}
\end{align}
Finally, let
\[
 \alpha:=\frac{2}{5+\eta},\qquad
 \theta:=\alpha\eta=\frac{2\eta}{5+\eta}=2-5\alpha.             \tag{4.5}
\]

\begin{theorem}[fixed-coarse power rate]\label{thm:rate}
Let \(r_*\) be the least integer satisfying
\[
 2^{\alpha r_*}\ge\max(1,2b).                                 \tag{4.6}
\]
For every \(r\ge r_*\),
\[
 |Q_{r,M}(q,p)-Q_{\infty,M}(q,p)|
 \le C_Q\,2^{-\theta r},                                      \tag{4.7}
\]
where
\[
 C_Q:=C_{\rm low}(1+b^{-1})^5+\frac{2^\eta C_*}{\pi\eta}.      \tag{4.8}
\]
Consequently,
\[
 \left|
 \omega_{M+r,L}^{(0)}(\alpha_M^{M+r}(W_M(\eps_Mq+ip)))
 -\omega_L^{(0)}(\beta_M(W_M(\eps_Mq+ip)))
 \right|
 \le\frac{C_Q}{4}\,2^{-\theta r}.                             \tag{4.9}
\]
\end{theorem}

\subsection*{Proof in Lamport form}

\LS{1}{1}{Choose an integer momentum split.}

For \(r\ge r_*\), put
\[
 J_r:=\left\lfloor b^{-1}2^{\alpha r}\right\rfloor.            \tag{4.10}
\]
Then \(J_r\ge1\),
\[
 bJ_r\le2^{\alpha r}\le2^r,\qquad
 bJ_r\ge2^{\alpha r}-b\ge\frac12\,2^{\alpha r},                \tag{4.11}
\]
and every mode \(k_n=\kappa n\) with \(|n|\le J_r\) lies in
\(\Gamma_{M+r}\).  The first inequality in (4.11) also permits use of
Lemma \ref{lem:filtertail}.

\LS{1}{2}{Bound the low-momentum part.}

For \(|n|\le J_r\), split the momentum integrand as
\begin{align}
 |\gamma_N|P_r|^2-\gamma|P_\infty|^2|
 &\le|\gamma_N-\gamma||P_r|^2
 +\gamma\bigl||P_r|^2-|P_\infty|^2\bigr|,                     \tag{4.12}\\
 |\gamma_N^{-1}|P_r|^2-\gamma^{-1}|P_\infty|^2|
 &\le|\gamma_N^{-1}-\gamma^{-1}||P_r|^2
 +\gamma^{-1}\bigl||P_r|^2-|P_\infty|^2\bigr|.                \tag{4.13}
\end{align}
Here every mask argument is \(\ell=\eps_Mk_n\).  Lemmas
\ref{lem:filtertail} and \ref{lem:dispersion}, together with
\(\eps_{M+r}=2^{-r}\eps_M\), give
\begin{align}
 \text{\(q\)-integrand error}
 &\le4^{-r}\eps_M^2A_q^2
   \left(\frac{k_n^4}{24m^3}+\frac{c_0k_n^2}{m}\right),        \tag{4.14}\\
 \text{\(p\)-integrand error}
 &\le4^{-r}\eps_M^2A_p^2
   \left(\frac{k_n^4}{24m}+c_0\gamma(k_n)k_n^2\right).         \tag{4.15}
\end{align}
Because
\[
 2J+1\le3(1+J),\quad |k_n|\le\kappa(1+J),\quad
 \gamma(k_n)\le m+|k_n|\le(m+\kappa)(1+J),                    \tag{4.16}
\]
the definition (4.3) yields
\[
 E_{\rm low}(r)
 \le4^{-r}C_{\rm low}(1+J_r)^5
 \le C_{\rm low}(1+b^{-1})^5\,2^{-(2-5\alpha)r}.              \tag{4.17}
\]

\LS{1}{3}{Bound both high-momentum tails.}

Lemma \ref{lem:envelope} and
\[
 \gamma_{M+r}(k)\le\gamma(k)\le m+|k|
 \le(m+\eps_M^{-1})(1+\eps_M|k|)                              \tag{4.18}
\]
show that the sum of the absolute fine and continuum integrands at
\(k_n=\kappa n\) is at most
\[
 C_*(1+b|n|)^{-1-\eta}.                                       \tag{4.19}
\]
The fine term is declared zero outside \(\Gamma_{M+r}\), exactly as in (1.4).
Since \(x\mapsto(1+bx)^{-1-\eta}\) is decreasing,
\[
 \sum_{|n|>J}(1+b|n|)^{-1-\eta}
 \le2\int_J^\infty(1+bx)^{-1-\eta}\dd x
 =\frac{2}{b\eta}(1+bJ)^{-\eta}.                              \tag{4.20}
\]
Multiplication by the covariance weight \(\eps_M/(2L)\), followed by
\(b=\pi\eps_M/L\) and (4.11), gives
\[
 E_{\rm high}(r)
 \le\frac{C_*}{\pi\eta}(1+bJ_r)^{-\eta}
 \le\frac{2^\eta C_*}{\pi\eta}\,2^{-\alpha\eta r}.             \tag{4.21}
\]

\LS{1}{4}{Match the exponents.}

By (4.5),
\[
 2-5\alpha=\alpha\eta=\theta.                                 \tag{4.22}
\]
Adding (4.17) and (4.21) proves (4.7)--(4.8).

\LS{1}{5}{Pass from covariance forms to states.}

For \(x,y\ge0\), the mean-value theorem gives
\[
 |e^{-x/4}-e^{-y/4}|
 =\frac14e^{-\zeta/4}|x-y|
 \le\frac14|x-y|                                               \tag{4.23}
\]
for some \(\zeta\) between \(x\) and \(y\).  Apply (4.23) to
\(Q_{r,M}\) and \(Q_{\infty,M}\) and use (4.7).  This is (4.9).
\QedStep

\begin{corollary}[lattice-spacing form]
Because \(\eps_{M+r}/\eps_M=2^{-r}\), (4.9) is equivalently
\[
 |\omega_{M+r,L}^{(0)}\circ\alpha_M^{M+r}(W_M(\xi))
 -\omega_L^{(0)}\circ\beta_M(W_M(\xi))|
 \le\frac{C_Q}{4}\left(\frac{\eps_{M+r}}{\eps_M}\right)^\theta.
                                                                    \tag{4.24}
\]
\end{corollary}

\section{Claim ledger}

\begin{longtable}{P{0.31\textwidth}P{0.17\textwidth}P{0.43\textwidth}}
\toprule
Claim & Status & Exact reason\\
\midrule
\endhead
Fixed-coarse free convergence & \Pass & The covariance formulas are (1.4) and
(1.8), and Theorem \ref{thm:rate} supplies a power rate.\\
Uniform finite \(D4\) filter tail & \Pass & Lemmas \ref{lem:envelope} and
\ref{lem:filtertail} prove both the global envelope and the low-frequency truncation
error from the exact four mask coefficients.\\
Finest-scale local state norm & \Fail & This theorem fixes \(M\) and sends
\(r=N-M\to\infty\); it does not include \(M=N\).\\
Interacting covariance reduction & \Fail & Interacting ground states are not Gaussian,
so no quadratic covariance symbol determines their Weyl expectations.\\
\bottomrule
\end{longtable}

\begin{thebibliography}{9}
\bibitem{MMST}
V.~Morinelli, G.~Morsella, A.~Stottmeister, Y.~Tanimoto,
\emph{Scaling limits of lattice quantum fields by wavelets},
Commun. Math. Phys. \textbf{387} (2021), 1299--1367.
\end{thebibliography}

\end{document}
